\documentclass[11pt,a4paper]{article}
\usepackage[utf8]{inputenc}
\usepackage[T1]{fontenc}
\usepackage[spanish,es-nodecimaldot]{babel}
\usepackage{amsmath,amssymb}
\usepackage[margin=2.5cm]{geometry}
\usepackage{enumitem}
\setlist[enumerate,1]{label=\textbf{\arabic*.},leftmargin=*,itemsep=1em}
\setlist[enumerate,2]{label=\alph*),leftmargin=1cm}
\setlength{\parindent}{0pt}
\title{\textbf{Solucionario de modelamiento\\Examen 01}}
\author{}
\date{}
\begin{document}
\maketitle
El orden de las variables sigue el orden de los conceptos o tipos en cada enunciado. Se usa $\mathbf{x}=(x,y,z)^T$.
\begin{enumerate}
\item \textbf{Servicios de un taller.} En orden, $x,y,z$ representan: costos de diagnóstico, ajuste y limpieza, en dólares.
\[\begin{cases}x+y+z=60\\y+z=35\\x-z=10\end{cases}\qquad A=\begin{pmatrix}1&1&1\\0&1&1\\1&0&-1\end{pmatrix},\quad b=\begin{pmatrix}60\\35\\10\end{pmatrix}.\]
\[\det(A)=-1,\qquad A^{-1}=\frac{1}{-1}\begin{pmatrix}-1&1&0\\1&-2&-1\\-1&1&1\end{pmatrix},\qquad \mathbf{x}=A^{-1}b=\begin{pmatrix}25\\20\\15\end{pmatrix}.\]
Por tanto, en el orden definido, los valores son 25, 20, 15 dólares. Como $\det(A)\ne0$, la solución es única.
\item \textbf{Alquiler de equipos.} En orden, $x,y,z$ representan: tarifas de alquiler por día de un proyector, una pantalla y un sistema de sonido, en dólares.
\[\begin{cases}x+y+z=65\\2x+y=80\\y-z=5\end{cases}\qquad A=\begin{pmatrix}1&1&1\\2&1&0\\0&1&-1\end{pmatrix},\quad b=\begin{pmatrix}65\\80\\5\end{pmatrix}.\]
Al sustituir la columna $i$ de $A$ por $b$, se obtiene $A_i$. Entonces
\[\det(A)=3,\quad \det(A_1)=90,\quad \det(A_2)=60,\quad \det(A_3)=45.\]
\[(x,y,z)=\left(\frac{90}{3},\frac{60}{3},\frac{45}{3}\right)=(30,20,15).\]
Por tanto, en el orden definido, los valores son 30, 20, 15 dólares por día. Como $\det(A)\ne0$, la solución es única.
\item \textbf{Kits de laboratorio.} En orden, $x,y,z$ representan: cantidades de kits de los tipos A, B y C.
\[\begin{cases}x+2y+z=11\\2x+y+z=12\\x+y+2z=13\end{cases}\qquad A=\begin{pmatrix}1&2&1\\2&1&1\\1&1&2\end{pmatrix},\quad b=\begin{pmatrix}11\\12\\13\end{pmatrix}.\]
Al sustituir la columna $i$ de $A$ por $b$, se obtiene $A_i$. Entonces
\[\det(A)=-4,\quad \det(A_1)=-12,\quad \det(A_2)=-8,\quad \det(A_3)=-16.\]
\[(x,y,z)=\left(\frac{-12}{-4},\frac{-8}{-4},\frac{-16}{-4}\right)=(3,2,4).\]
Por tanto, en el orden definido, los valores son 3, 2, 4 kits. Como $\det(A)\ne0$, la solución es única.
\newpage
\item \textbf{Mezcla de concentrados.} En orden, $x,y,z$ representan: volúmenes de los concentrados A, B y C, en litros.
\[\begin{cases}x+y+z=9\\2x+y=7\\y+3z=15\end{cases}\qquad A=\begin{pmatrix}1&1&1\\2&1&0\\0&1&3\end{pmatrix},\quad b=\begin{pmatrix}9\\7\\15\end{pmatrix}.\]
\[\det(A)=-1,\qquad A^{-1}=\frac{1}{-1}\begin{pmatrix}3&-2&-1\\-6&3&2\\2&-1&-1\end{pmatrix},\qquad \mathbf{x}=A^{-1}b=\begin{pmatrix}2\\3\\4\end{pmatrix}.\]
Por tanto, en el orden definido, los valores son 2, 3, 4 litros. Como $\det(A)\ne0$, la solución es única.
\item \textbf{Producción de paneles.} En orden, $x,y,z$ representan: cantidades de paneles de los modelos A, B y C.
\[\begin{cases}2x+y+z=13\\x+2y=10\\x+z=6\end{cases}\qquad A=\begin{pmatrix}2&1&1\\1&2&0\\1&0&1\end{pmatrix},\quad b=\begin{pmatrix}13\\10\\6\end{pmatrix}.\]
Al sustituir la columna $i$ de $A$ por $b$, se obtiene $A_i$. Entonces
\[\det(A)=1,\quad \det(A_1)=4,\quad \det(A_2)=3,\quad \det(A_3)=2.\]
\[(x,y,z)=\left(\frac{4}{1},\frac{3}{1},\frac{2}{1}\right)=(4,3,2).\]
Por tanto, en el orden definido, los valores son 4, 3, 2 paneles. Como $\det(A)\ne0$, la solución es única.
\item \textbf{Entradas a una exposición.} En orden, $x,y,z$ representan: precios de las entradas general, estudiantil e infantil, en dólares.
\[\begin{cases}x+y+z=22\\x-2y=0\\2x+z=28\end{cases}\qquad A=\begin{pmatrix}1&1&1\\1&-2&0\\2&0&1\end{pmatrix},\quad b=\begin{pmatrix}22\\0\\28\end{pmatrix}.\]
\[\det(A)=1,\qquad A^{-1}=\frac{1}{1}\begin{pmatrix}-2&-1&2\\-1&-1&1\\4&2&-3\end{pmatrix},\qquad \mathbf{x}=A^{-1}b=\begin{pmatrix}12\\6\\4\end{pmatrix}.\]
Por tanto, en el orden definido, los valores son 12, 6, 4 dólares. Como $\det(A)\ne0$, la solución es única.
\newpage
\item \textbf{Transporte de carga.} En orden, $x,y,z$ representan: cantidades de vehículos de los tipos A, B y C.
\[\begin{cases}x+y+z=9\\x+2y+3z=17\\2x+y+z=12\end{cases}\qquad A=\begin{pmatrix}1&1&1\\1&2&3\\2&1&1\end{pmatrix},\quad b=\begin{pmatrix}9\\17\\12\end{pmatrix}.\]
Al sustituir la columna $i$ de $A$ por $b$, se obtiene $A_i$. Entonces
\[\det(A)=1,\quad \det(A_1)=3,\quad \det(A_2)=4,\quad \det(A_3)=2.\]
\[(x,y,z)=\left(\frac{3}{1},\frac{4}{1},\frac{2}{1}\right)=(3,4,2).\]
Por tanto, en el orden definido, los valores son 3, 4, 2 vehículos. Como $\det(A)\ne0$, la solución es única.
\item \textbf{Calibración de señales.} En orden, $x,y,z$ representan: intensidades de las fuentes A, B y C, en unidades de intensidad.
\[\begin{cases}x+2y=8\\y+z=4\\2x+y+z=8\end{cases}\qquad A=\begin{pmatrix}1&2&0\\0&1&1\\2&1&1\end{pmatrix},\quad b=\begin{pmatrix}8\\4\\8\end{pmatrix}.\]
\[\det(A)=4,\qquad A^{-1}=\frac{1}{4}\begin{pmatrix}0&-2&2\\2&1&-1\\-2&3&1\end{pmatrix},\qquad \mathbf{x}=A^{-1}b=\begin{pmatrix}2\\3\\1\end{pmatrix}.\]
Por tanto, en el orden definido, los valores son 2, 3, 1 unidades de intensidad. Como $\det(A)\ne0$, la solución es única.
\item \textbf{Paquetes de impresión.} En orden, $x,y,z$ representan: cantidades de paquetes de los tipos A, B y C.
\[\begin{cases}x+y+z=6\\2x+y+3z=10\\x+2y+z=9\end{cases}\qquad A=\begin{pmatrix}1&1&1\\2&1&3\\1&2&1\end{pmatrix},\quad b=\begin{pmatrix}6\\10\\9\end{pmatrix}.\]
Al sustituir la columna $i$ de $A$ por $b$, se obtiene $A_i$. Entonces
\[\det(A)=-1,\quad \det(A_1)=-2,\quad \det(A_2)=-3,\quad \det(A_3)=-1.\]
\[(x,y,z)=\left(\frac{-2}{-1},\frac{-3}{-1},\frac{-1}{-1}\right)=(2,3,1).\]
Por tanto, en el orden definido, los valores son 2, 3, 1 paquetes. Como $\det(A)\ne0$, la solución es única.
\end{enumerate}
\end{document}
